\section{Exact lifts and local conventions}\label{sec:foundations}
Throughout, cone-dimension and matrix-order caps are integers.

\subsection{Exact lifts and genuine certificates}
A \emph{proper cone} is a closed, convex, pointed cone with nonempty
interior in its finite-dimensional real ambient space. Its dual is
$K^*=\{y:\ip{x}{y}\geq0\text{ for all }x\in K\}$.
Let $C\subset\R^s$ be a compact convex body with nonempty interior and
$0\in\intr C$. Its polar body is
$C^\circ=\{v\in\R^s:\ip{x}{v}\leq1\text{ for every }x\in C\}$.
An exact lift is a triple
\[
 \cL=(K,A,\pi),\qquad C=\pi(K\cap A),
\]
where $K=\prod_{i=1}^kK_i$, $A$ is an affine subspace, and
$\pi:A\to\R^s$ is affine. We do not assume $K\cap A$ is compact.
A lift is \emph{proper} if $A\cap\intr K\ne\varnothing$.
For relative Slater feasibility we use the minimal-face reduction of the
following lemma.

\begin{lemma}[Reduction and certificate normalization]\label{lem:reduction}
Unrestricted free variables in a finite affine lift of a compact body can
be eliminated without changing its cone factors. The cone can then be
replaced by the smallest face containing all feasible points, within the
span of that face; the resulting lift is proper. For a proper lift and
every $v\in\partial C^\circ$, the set
\begin{equation}\label{eq:certificate-fiber}
 \cD_{\cL}(v)=\{Y\in K^*: \ip{X}{Y}=1-\ip{\pi(X)}{v}
                    \text{ for every }X\in A\}
\end{equation}
is nonempty. Thus certificates represent the support slack on all of
$A$, including points outside $K$.
\end{lemma}
\begin{proof}
Write the original affine equalities in cone coordinates $X$ and free
coordinates $z$. A free direction that preserves these equalities while
fixing $X$ must have zero output direction: otherwise its entire line
projects into the compact body. Consequently the output depends affinely
only on $X$, on the affine subspace of $X$ for which the equalities admit
some $z$. Linear elimination gives the claimed representation.

Choose $X_0\in\ri(K\cap A)$, and let $F$ be the minimal face of $K$
containing $X_0$. For every $X\in K\cap A$, the relative-interior property
gives $X_0-\epsilon(X-X_0)\in K\cap A$ for some $\epsilon>0$.
The face property therefore gives $X\in F$. Conversely $F\cap A\subset
K\cap A$. The point $X_0$ belongs to $\ri F$, so replacing $K$ by $F$
and $A$ by $A\cap\spanop F$ gives Slater feasibility in $\spanop F$.
Faces of a product are products of faces.

For the last assertion write $A=\{X:MX=b\}$ and extend the affine output
as $\pi(X)=TX+t$. Maximizing $\ip{v}{TX+t}$ over $K\cap A$ has value one.
Slater conic duality gives a multiplier $\lambda$ with
$Y=M^*\lambda-T^*v\in K^*$ and
$\ip{b}{\lambda}=1-\ip{t}{v}$. Hence for every $X\in A$,
$\ip{X}{Y}=\ip{b}{\lambda}-\ip{TX}{v}=1-\ip{\pi(X)}{v}$.
This is the duality step in the proper-lift direction of
\cite[Theorem~2.4]{GPT2013}.
\end{proof}

We call the elements of $\cD_{\cL}(v)$ the \emph{genuine certificates}
of $v$, and $\cD_{\cL}(v)$ its \emph{certificate fiber}.

For symmetric cones, every face is the cone of squares of a Peirce-one
subalgebra, and Jordan ranks in a face agree with ambient ranks.
Throughout, ranks, determinants, and barriers of a relative-Slater lift
are computed \emph{after} this reduction. Certificates then lie in the dual
of the reduced face and exist even if the original presentation fails
Slater's condition. A face of a dictionary factor need not belong
to the dictionary, but its capacities never exceed those of the factor.

The ball $C=\B_2^s=\{x:\norm{x}_2\leq1\}$ is its own polar, so its
supports are the unit vectors $v\in\Sph^{s-1}$. Define
\begin{equation}\label{eq:rank-invariant}
 r_{\cL}(v)=\min_{Y\in\cD_{\cL}(v)}\sum_i\jrank Y_i,
 \qquad R(\cL)=\max_{v\in\Sph^{s-1}}r_{\cL}(v).
\end{equation}
Both extrema exist because they are taken over nonempty subsets of a
finite set of integers; compactness of the certificate fiber is not
needed. No certificate is zero, because the support slack is not
identically zero on $C$; this also gives the lower bound one in the
interval case $s=1$, treated separately.

\subsection{Regularity obtained from the lift}
We call data \emph{definable} if they are definable, allowing real
parameters, in one fixed o-minimal expansion of the ordered real field.
Only the bound for general proper cones below needs this generality. All
symmetric-cone data and selections used for the rank theorem are
semialgebraic. Definable choice and finite differentiable stratification
are standard~\cite{vdDM1996}; in the semialgebraic setting see~\cite{BCR1998}.

\begin{lemma}[Generic regular selections]\label{lem:selection}
For a proper definable lift, primal lifts of boundary points and
certificates of supports $v\in\partial C^\circ$ can be chosen definably.
On a definable $C^2$ boundary patch, their compositions with the definable
normalized outward-normal contact map are $C^1$ in definable local coordinates on a dense relatively
open subset of the patch. For a symmetric-cone lift of a Euclidean ball,
certificates can also attain $r_{\cL}(v)$, and the common
regular subset can be chosen semialgebraic, with a lower-dimensional
complement in $\Sph^{s-1}$.
\end{lemma}
\begin{proof}
The primal fiber and the fiber~\eqref{eq:certificate-fiber} are nonempty
closed definable sets; the identity on $A$ amounts to finitely many affine
equalities. Their minimum-Euclidean-norm points are unique, since both
fibers are convex, and their graphs are definable by quantified formulas.
For minimum Jordan rank, each condition $\sum_i\jrank Y_i=j$ is
semialgebraic by spectral polynomial equations and inequalities.
Partition the support sphere according to the least realized $j$ and
apply semialgebraic choice in each resulting nonempty fiber. This choice
need not minimize norm and need not be globally continuous.

Finite $C^1$ cell decomposition, refined for all the maps at once, makes
them $C^1$ on each full-dimensional cell. The union of the interiors of
those cells is relatively open and dense; its complement has smaller
dimension. In the semialgebraic case this complement has surface measure
zero. A definable $C^2$ boundary patch and its normalized normal map are
eligible for the same decomposition. Faces used in reduction are definable:
for $X_0\in\ri F$,
$F=\{X\in K:\exists\epsilon>0, X_0-\epsilon X\in K\}$.
\end{proof}

\section{Local curvature and the size of cone factors}\label{sec:curvature}
A \emph{contact} is a pair $(x,v)$ with $x\in\partial C$,
$v\in\partial C^\circ$, and $\ip{x}{v}=1$.

Let $x:U\to\partial C$ be a definable $C^2$ local parametrization, where
$U\subset\R^{s-1}$, and let $n(u)$ be its outward unit normal. Set
\[
 y(u)=\frac{n(u)}{\ip{x(u)}{n(u)}}.
\]
The denominator is positive because $0\in\intr C$, and $y(u)$ lies on
$\partial C^\circ$. As in Lemma~\ref{lem:selection}, select a primal
lift $X(u)=(X_i(u))_i$ of $x(u)$ and a certificate
$Y(u)=(Y_i(u))_i\in\cD_{\cL}(y(u))$. The whole-slice identity gives
\begin{equation}\label{eq:local-slack}
 1-\ip{x(u)}{y(v)}=\sum_i\ip{X_i(u)}{Y_i(v)}.
\end{equation}
Each summand is nonnegative and vanishes on the diagonal. The negative of
the mixed second derivative at the diagonal point $u=v=0$ is the bilinear
form
\begin{equation}\label{eq:mixed-curvature}
 \mathsf H=Dx(0)^*Dy(0)
 =\frac{Dx(0)^*Dn(0)}{\ip{x(0)}{n(0)}}.
\end{equation}
At a strictly positively curved point it has rank $s-1$.

\begin{lemma}[Universal local bound]\label{lem:universal-channel}
Let $K\subset E$ be a proper cone, $m=\dim E$, and let
$X:U\to K$, $Y:U\to K^*$ be $C^1$ with
$\ip{X(u)}{Y(u)}=0$. The bilinear form
$-\ip{DX(0)[h]}{DY(0)[\ell]}$ has rank at most $m-2$ if
$X(0),Y(0)\ne0$, and rank zero if either is zero. In particular rays
contribute no mixed curvature.
\end{lemma}
\begin{proof}
Put $a=X(0)$ and $b=Y(0)$. When both are nonzero, nonnegativity of
$\ip{X(th)}b$ for positive and negative $t$ shows $DX(0)[h]\in b^\perp$.
Similarly $DY(0)[\ell]\in a^\perp$. The pairing from $b^\perp$ to
$(a^\perp)^*$ has kernel
$b^\perp\cap\spanop\{a\}=\spanop\{a\}$, since $\ip ab=0$.
Its rank is $(m-1)-1=m-2$. Pullback by the derivatives cannot increase
rank. If $a=0$, a two-sided derivative of $X$ lies in $K\cap(-K)=\{0\}$;
the case $b=0$ is identical.
\end{proof}

\begin{theorem}[Curvature bounds]\label{thm:curvature}
Suppose a compact $s$-dimensional body has a relatively open $C^2$
boundary patch with strictly positive curvature at some point, and an
exact lift by a finite product of definable proper cones of dimensions
$m_i$. Then
\begin{equation}\label{eq:universal-capacity}
 s-1\leq\sum_i\max\{m_i-2,0\}.
\end{equation}
If the factors are simple symmetric cones of ranks $r_i$ and Peirce
constants $a_i$, one can choose a contact at which the complementary
primal and certificate Jordan ranks $p_i,q_i$ satisfy
\begin{equation}\label{eq:jordan-capacity}
 s-1\leq\sum_i a_ip_iq_i
 \leq\sum_i a_i\left\lfloor\frac{r_i^2}{4}\right\rfloor.
\end{equation}
A ray has $r_i=1$, $a_i=0$ in this notation.
\end{theorem}

We first prove the sharper local bound for symmetric cones. In a Euclidean Jordan algebra $V$, write $x\circ y$ for
its product, $e$ for its identity, and $\ip xy=\tr(x\circ y)$ for the
trace inner product. The cone $K$ is the cone of squares, and $x\succeq0$
means $x\in K$. A Jordan frame $c_1,\ldots,c_r$ gives the orthogonal
Peirce decomposition $V=\bigoplus_{i\leq j}V_{ij}$, where
$V_{ii}=\R c_i$ and $\dim V_{ij}=a$ for $i<j$ in a simple algebra.
The support idempotent of $x\succeq0$ is the sum of the frame idempotents
with positive eigenvalues; their number is $\jrank x$.

\begin{lemma}[Mixed Peirce channels]\label{lem:peirce-channel}
Let $X,Y:U\to K$ be $C^1$ and complementary:
$\ip{X(u)}{Y(u)}=0$. If $X(0)$ and $Y(0)$ have ranks $p$ and $q$,
then $p+q\leq r$, and
$-\ip{DX(0)[h]}{DY(0)[\ell]}$ is a positive semidefinite symmetric
bilinear form of rank at most $apq$. Neither locally constant ranks nor
strict complementarity is required.
\end{lemma}
\begin{proof}
Positive orthogonal elements are Jordan complementary and diagonalize in
a common Jordan frame~\cite[Chapters~III--IV]{FK1994}. Thus
\[
 X(0)=\sum_{i\in I}\lambda_i c_i,\quad
 Y(0)=\sum_{j\in J}\mu_j c_j,
 \qquad \lambda_i,\mu_j>0,\quad I\cap J=\varnothing.
\]
For every positive $z$ in the Peirce subalgebra on the complementary
indices $I^c$, $\ip{X(th)}z\geq0$ vanishes at $t=0$, so its derivative
there is zero. Since the positive cone of that subalgebra spans it, the
derivative of $X$ has zero projection onto it. Likewise the derivative of $Y$ vanishes
on the $J^c$ subalgebra. The only possible overlap of the two derivatives
is therefore $\bigoplus_{i\in I,j\in J}V_{ij}$, of dimension $apq$.

Differentiate $X(u)\circ Y(u)=0$. Its $V_{ij}$ component is
$\mu_j(DX[h])_{ij}+\lambda_i(DY[h])_{ij}=0$, because multiplication
by either endpoint idempotent acts by one half on $V_{ij}$. Hence
\begin{equation}\label{eq:peirce-gram}
 -\ip{DX[h]}{DY[\ell]}
 =\sum_{i\in I,j\in J}\frac{\mu_j}{\lambda_i}
       \ip{(DX[h])_{ij}}{(DX[\ell])_{ij}}.
\end{equation}
This is the required Gram form. If $I$ or $J$ is empty, the form is zero,
as also follows from the pointed-vertex argument of
Lemma~\ref{lem:universal-channel}.
\end{proof}

\begin{proof}[Proof of Theorem~\ref{thm:curvature}]
Reduce the lift by Lemma~\ref{lem:reduction}. Curvature is strictly
positive on an open subpatch. There choose a contact in
the common regular set given by Lemma~\ref{lem:selection}. Differentiating~\eqref{eq:local-slack} once
in each independent variable and using rank subadditivity with
Lemma~\ref{lem:universal-channel} gives~\eqref{eq:universal-capacity}.
A reduced face has dimension no larger than its original factor, so the
original dimensions give valid bounds. For symmetric cones apply
Lemma~\ref{lem:peirce-channel} instead. The inequality
$pq\leq\lfloor r^2/4\rfloor$ follows from $p+q\leq r$.
Faces of a simple symmetric cone have the same Peirce constant at smaller
rank, except that rank-one faces are rays. For the Albert cone its
rank-two faces are spin factors with Peirce constant eight. Their
capacities $a\lfloor r^2/4\rfloor$ likewise cannot exceed that of the
original factor.
\end{proof}

The bound $apq$ is attained locally. The inner derivations
$4[L_z,L_{c_i}]$, $z\in V_{ij}$, independently rotate the support
directions $i$ and $j$, and their exponentials are Jordan automorphisms.
Applied simultaneously to $X(0),Y(0)$, their derivatives in $V_{ij}$ are
$\lambda_i z,-\mu_jz$. Thus all $apq$ channels in \eqref{eq:peirce-gram}
can be realized. This local statement does not claim that a specified
collection of channels integrates into an exact lift.

Write $Q_m=\{(t,z)\in\R\times\R^{m-1}:t\geq\norm{z}_2\}$ for
the Lorentz cone of total cone-space dimension $m$.

\begin{corollary}[Least factor count and dimension for a ball under a dimension cap]\label{cor:dimension-resources}
Fix $s\geq2$ and an integer $d\geq3$. Among exact lifts of $\B_2^s$ by
arbitrary definable proper cone factors of dimension at most $d$, the
minimum number of factors and the minimum total cone-space dimension are
\[
 k_{\min}=\left\lceil\frac{s-1}{d-2}\right\rceil,
 \qquad M_{\min}=s-1+2k_{\min}.
\]
Both minima are attained by the same Lorentz norm tree, constructed in
the proof. Among
symmetric cone lifts under this dimension cap, the least total ambient
Jordan rank is $2k_{\min}$, again attained by that lift.
\end{corollary}
\begin{proof}
Let $k_+$ count factors of dimension at least two and $k_1$ count rays.
Theorem~\ref{thm:curvature} implies
$s-1\leq k_+(d-2)$ and
$M\geq s-1+2k_++k_1$, proving the first two lower bounds.
For symmetric factors only non-ray irreducible factors can contribute
curvature, and each has rank at least two. Decomposing reducible factors
cannot increase their dimensions or change total rank, so total rank is
at least $2k_{\min}$.

For attainment we use a variable-arity form of the recursive norm
representation in~\cite[Section~2, equation~(5)]{BTN2001}.
Put $k=k_{\min}$ and write $s-1=b_1+\cdots+b_k$, with
$1\leq b_i\leq d-2$. Form a rooted tree whose internal nodes are a chain.
Node $i<k$ has $b_i$ leaf children and the next internal node as one
additional child; the last node has $b_k+1$ leaf children. It has $s$
leaves. Assign the coordinates of $x\in\R^s$ to those leaves and an
auxiliary scalar $t_i$ to each internal node, and impose
\[
 t_i\geq\norm{\text{child values of node }i}_2
 \quad\text{in }Q_{b_i+2},\qquad t_1=1.
\]
Recursive elimination of the scalars gives exactly
$\norm{x}_2\leq1$. Strict feasibility
holds at $x=0$ by choosing the nonroot scalars positive and successively
smaller. The construction uses $k$ rank-two factors and total dimension
$\sum_i(b_i+2)=s-1+2k$, as required.
\end{proof}

The total ambient Jordan rank is the parameter of the standard product
log-determinant barrier on the ambient cone and remains valid after
restriction to the affine slice. It need not be the least parameter of
the restricted barrier, and it gives no lower bound for a barrier defined
intrinsically on the projected ball. These distinctions matter when
curvature bounds are used to compare formulations.
